\documentclass[11pt]{article}
\usepackage[T1]{fontenc}
\usepackage{lmodern}
\usepackage[margin=1in]{geometry}
\usepackage{amsmath,amssymb,amsthm,mathtools}
\usepackage{booktabs,array,longtable,enumitem}
\usepackage{microtype,xcolor,hyperref,fancyhdr,listings}
\hypersetup{colorlinks=true,linkcolor=blue!45!black,urlcolor=blue!50!black,citecolor=blue!45!black,pdftitle={Fixed-input Diophantine forms of ZFC consistency},pdfauthor={Research collaboration for ProveIt}}
\pagestyle{fancy}\fancyhf{}\fancyhead[L]{\small Fixed-input Diophantine forms of ZFC consistency}\fancyhead[R]{\small ProveIt}\fancyfoot[C]{\thepage}
\setlength{\headheight}{14pt}
\setlength{\emergencystretch}{2em}
\newcommand{\NN}{\mathbb N}
\newcommand{\ZZ}{\mathbb Z}
\newcommand{\QQ}{\mathbb Q}
\newcommand{\Pos}{\mathbb Z_{>0}}
\newcommand{\Con}{\operatorname{Con}}
\newcommand{\Sol}{\operatorname{Sol}}
\newcommand{\pc}{\operatorname{pc}}
\newtheorem{theorem}{Theorem}[section]
\newtheorem{lemma}[theorem]{Lemma}
\newtheorem{proposition}[theorem]{Proposition}
\theoremstyle{remark}\newtheorem{remark}[theorem]{Remark}
\lstset{basicstyle=\ttfamily\footnotesize,breaklines=true,columns=fullflexible,keepspaces=true,frame=single,rulecolor=\color{black!20}}
\title{\textbf{Fixed-input Diophantine forms of ZFC consistency}\\[5pt]\large Seven equations in sixteen positive unknowns;\\single equations in fifteen positive unknowns}
\author{A research contribution for Vladimir Reshetnikov's ProveIt repository}
\date{22 October 2026}
\begin{document}
\maketitle
\begin{abstract}
We construct a fixed system $\mathcal S_{16}$ of seven polynomial equations in sixteen strictly positive integer unknowns, and fixed polynomials $\Phi_{15}$ and $\Psi_{15}$ in fifteen such unknowns, with
\[
\mathrm{ZFC}\vdash\Con(\mathrm{ZFC})\ \longleftrightarrow\
\neg\Sol(\mathcal S_{16})\ \longleftrightarrow\
\neg\Sol(\Phi_{15}=0)\ \longleftrightarrow\
\neg\Sol(\Psi_{15}=0).
\]
The system has maximum expanded degree $34$ and a shared evaluation circuit of $66$ arithmetic operations. The two fifteen-variable polynomials have respective exact degrees $295$ and $346$, with $82$ and $80$ operations. These counts include every variable-dependent operation in the supplied circuits; fixed integer coefficients are counted as constants.

The main new reduction specializes the repository's universal polynomial to one fixed input and removes its entire Pell input bridge. A constructive inverse restores the deleted positive witnesses, so removing the bridge introduces no additional solutions. A quadratic norm then removes the strong Pell root. The coefficients are fully specified finite sums derived from the literal K215 machine, not an unspecified recursively enumerable index. They are astronomically large and are supplied by exact definitions, not expanded digit strings. The contribution is an ordinary mathematical proof with reproducible arithmetic and finite-machine checks; it is not a new Lean or Rocq verification of the complete Diophantine reduction, and it makes no global minimality claim.
\end{abstract}
\newpage
\tableofcontents
\newpage

\section{Result, domains, and meaning of simplicity}

Throughout, $\NN=\{0,1,2,\ldots\}$. Our primary domain is $\Pos$. Every occurrence of a computed quantity below is a polynomial abbreviation in the listed unknowns. These abbreviations do not introduce further existential variables.

\begin{theorem}[Concrete Diophantine consistency forms]\label{thm:main}
Use the fixed integers specified in Section~\ref{sec:coefficients}, and the explicit equations in Sections~\ref{sec:system} and~\ref{sec:fifteen}. Then ZFC proves that each of the following is equivalent to the existence of a finite syntactic ZFC refutation:
\begin{enumerate}[label=(\roman*),nosep]
\item $\mathcal S_{16}$ has a solution in $\Pos^{16}$;
\item $\Phi_{15}=0$ has a solution in $\Pos^{15}$;
\item $\Psi_{15}=0$ has a solution in $\Pos^{15}$.
\end{enumerate}
Consequently their lack of solutions is provably equivalent in ZFC to $\Con(\mathrm{ZFC})$.
\end{theorem}

The proof uses the accepted K215 calibration and the existing native-tableau/Pell compiler theorems at a fixed repository commit. The new parts are the literal machine normalization and coefficient specification, the fixed-input elimination with a positive inverse, and their composition with the fifteen-variable norm elimination. Dependencies and the distinction between mathematical and executable evidence are recorded in Section~\ref{sec:evidence}.

There is no uniquely defined ``simplest equation'' without a cost model. Number of unknowns, number of equations, degree, coefficient height, and expression length lead to different choices. We recommend $\mathcal S_{16}$ for a readable mathematical presentation and $\Phi_{15}$ for a single equation with fewer unknowns. $\Psi_{15}$ saves two arithmetic operations at the cost of higher degree.

\begin{center}
\begin{tabular}{lrrrr}
\toprule
Construction & Unknowns & Equations & Degree & Operations\\
\midrule
$\mathcal S_{16}$ &16&7&34&66\\
Sum of its squared residuals &16&1&68&86\\
Scaled six-factor form &16&1&157&74\\
$\Phi_{15}$, local norm guard &15&1&295&82\\
$\Psi_{15}$, outer norm &15&1&346&80\\
\bottomrule
\end{tabular}
\end{center}

For a system, the degree is the largest expanded equation degree. Operations are binary additions, subtractions, and multiplications, including multiplication by fixed numerals. A system's equality comparisons are listed separately; the sum-of-squares and single-polynomial rows pay their entire finalizers. The numbers above do not measure the bit complexity of expanding the coefficients or evaluating the enormous integers in an actual solution.

Classical universal constructions can use fewer unknowns: the pairs $(9,1.64\cdot10^{45})$ and $(58,4)$ are universal over nonnegative integers, as recalled by Bayer and David~\cite{BayerDavid2025} from Jones~\cite{Jones1982}. Shifting the unknowns transfers these bounds to positive integers. Thus fifteen is not an absolute variable minimum. Our objective is an explicit, compact, repository-bound construction with a clear proof and a fully emitted arithmetic circuit. No lower bound over all Diophantine representations is asserted.

\section{Pinned source and the exact computation}

All repository references in this contribution are pinned to
\begin{center}\ttfamily\small cdd582bec13f908df3711fa9e05d1ddac368a5db.\end{center}
The accepted binary K215 machine has $215$ working states, a separate halt, a two-way tape, blank symbol $0$, initial head coordinate $0$, and initial state $0$. Its dense table has SHA-256
\begin{center}\ttfamily\scriptsize fdbcef0825886035c3091c8537a69535bc96f9e265f0ae45863543d40ae70090.\end{center}
The native, sparsely numbered presentation has SHA-256
\begin{center}\ttfamily\scriptsize 9683d6ed9628e6ce407c315b08f5da3d063da94e92706f275ca0f33e8272e4b2.\end{center}
The supplied checker verifies their exact state-renaming correspondence, including all $430$ rows, start state, and halt. Our normalizer reads the native table. It therefore simulates the same literal machine as the dense formal endpoint, with a different numbering.

The accepted Lean endpoint is
\begin{lstlisting}
KEndpoint215.blank_halts_iff_zfc_refutation :
  BlankHalts machine215 <-> BProv ZFCax_s [] Form.fBot
\end{lstlisting}
and the accepted Rocq endpoint is
\begin{lstlisting}
K215F.CorrectDense215.tm215_halts_iff_ZFC_inconsistent :
  Dense.halts dense_tm Dense.c0 <-> StandardZFC |-TC bottom
\end{lstlisting}
The second display uses ASCII typography for the final judgment. We use their ordinary mathematical content: blank halting is equivalent to a finite ZFC derivation of falsity. The repository records successful independent endpoint checks, with the Lean classical infrastructure axioms only and closed Rocq endpoint assumptions~\cite{K215Verification,K215Candidate}. Those historical records were inspected; a fresh full endpoint build is not claimed here.

\subsection{A completely specified one-sided normalizer}

To feed the native tableau compiler, normalize the machine to a tape with coordinates at most $0$. Physical cell $-j$ stores the two original tape bits at $j$ and $-j-1$. Finite control stores the active track. On the nonnegative track, an original right move becomes a physical left move; on the negative track, the direction is preserved. A move between original cells $0$ and $-1$ switches tracks at physical coordinate $0$ without moving. An origin flag is preserved on that physical cell.

The ten tape symbols are blank $0$, unary input $1$, temporary endpoint $2$, three further ordinary paired symbols, and four origin-marked paired symbols. Six control states precede the $2\cdot215$ simulation states: the initial state, the required first-step state, erase, return, halt, and a rejecting trap. Hence the normalized table has $436$ states, ten symbols, and $4360$ total transition rows. Its halt and trap are stationary loops; only the former is accepting.

The forced first transition writes $2$ and stays. The next transition moves left, erases the finite unary run, returns right to $2$, changes it to the marked zero pair, and enters the simulation of native state $0$. For unary length $n\ge1$, this prelude takes exactly $2n+2$ transitions. It never visits a positive physical coordinate. A native halting transition performs the specified write and move before entering the normalized halt.

The later raw81/raw76 tableau convention uses this stationary first transition; the earlier unary-tableau note used a left-moving convention. The supplied Start/End windows use the later convention. They also use the raw76 and modified75 ordering, Start $=0$ and End $=1$; the older raw81 ordering must not be substituted.

The machine ignores the unary input after erasing it. It accepts every nonempty unary input if and only if K215 halts on blank tape. We need only the particular input selected by the compiler's marker distance $2$: three initialized $I$ cells followed by the initial head cell, namely four unary ones. There is no claim that our machine rejects odd inputs as a general-purpose doubled-input wrapper would. Behavior on unqueried inputs has no role in this fixed-instance proof.

\section{The fixed coefficients}\label{sec:coefficients}

The main equations use six fixed integer ports:
\[
 B-1,\quad K_0,\quad M_C,\quad M_F,\quad W_0,\quad W_0+2d.
\]
Here $M_F$ is the \emph{shifted} mask. Omitting its shift changes the polynomial and invalidates the compiler interface.

\subsection{Finite windows, with no unspecified program index}

Use the helical marked-tableau alphabet for the normalized machine. There are two boundary symbols, $4370=10(436+1)$ ordinary tape/head payload symbols, and four first-row phase symbols $L,I,Q,R$. Thus the base alphabet has
\[
 a_*=4376
\]
symbols. Their order is $V,H$, then ordinary payloads in tape-symbol order and head-state order (no head first), then $L,I,Q,R$. The literal state and tape-symbol numberings are in the emitted machine table. The local predicate is stated explicitly in the supplied \texttt{fixed\_machine.py}: vertical boundary propagation, horizontal propagation within a strip, nonadjacent parallel boundaries, the first-row phase graph, exact deterministic three-neighbor updates, no head escape, and a halted head on the last row. It tests the five-cell cross of a $3\times3$ window; its four corners are unrestricted alphabet symbols.

List the allowed windows as $\mathcal W_0,\ldots,\mathcal W_{k_*-1}$, with the prescribed Start window first, the prescribed End window second, and all remaining windows in lexicographic base-alphabet order. The exact count, computed by factoring the incoming-head alternatives in the cross, is
\[
 k_*=115264654755851622862362693632.
\]
This count and the canonical finite predicate completely determine the list. Enumerating all its members is unnecessary and infeasible in practice. It is a fixed finite list independent of halting: membership in it is a finite local calculation, not an acceptance query.

\subsection{Explicit geometry and finite-sum definitions}

The following definitions make the large coefficients concrete. They are also implemented by a terminating, though infeasible, materializer in \texttt{coefficient\_recipe.py}. Let
\[
 A_*=2^{100},\qquad b=L_*=5^{45},\qquad d=5^{90},
 \qquad R_*=2^b,\qquad B=2^d.
\]
Put
\begin{align*}
 j_*&=57632327377925811431181340154,\\
 e_*&=k_*+12a_*,& M_*&=k_*+15a_*+1,\\
 H_*&=51M_*+3a_*+1,\\
 T_1&=H_*+54M_*+a_*+1,&T_2&=T_1+54M_*+1.
\end{align*}
For $r,c\in\{0,1,2\}$ and $0\le s<a_*$ define
\[
 p(r,c,s)=k_*+3a_*+(2-r)3a_*+(2-c)a_*+s,
 \qquad j(r,c,s)=1+(3r+c)a_*+s.
\]
The four anchors are $M_*,3M_*,9M_*,27M_*$. The native support is
\[
 \mathcal E=\{0,\ldots,k_*-1\}\cup\{p(r,c,s)\}\cup
 \{e_*\}\cup\{M_*,3M_*,9M_*,27M_*\}.
\]
Define coefficients $c_e$ by
\begin{align*}
 c_i&=1+\sum_{r,c=0}^2 A_*^{1+(3r+c)a_*+\mathcal W_i[3r+c]}
             +\sum_{\ell=1}^4 A_*^{9a_*+\ell}\quad(0\le i<k_*),\\
 c_{p(r,c,s)}&=A_*^{j(r,c,s)},\\
 c_{(1,3,9,27)_\ell M_*}&=A_*^{9a_*+\ell}\quad(1\le\ell\le4),
\end{align*}
and $c_e=0$ elsewhere. The indicated support parts are disjoint. Set
\[
 \mu_*=A_*-2+\sum_{j=1}^{j_*} A_*^j.
\]
All sums in the next display have explicit, finite limits:
\begin{align*}
 D_C={}&R_*^{3a_*}+R_*^{H_*+a_*}+R_*^{8M_*}+R_*^{24M_*}
       +\sum_e c_e\bigl(R_*^{T_1-e}+R_*^{T_2-e}\bigr),\\
 D_R={}&R_*^{H_*},\qquad K_0=D_C+BD_R,\\
 M_C={}&B-1-\sum_{e\in\mathcal E\setminus\{1\}}R_*^e-2R_*^{e_*},\\
 M_F^{\rm nat}={}&\mu_*\bigl(R_*^{T_1}+R_*^{T_2}\bigr)
       +\sum_{\substack{0\le r\le1,\ 0\le c\le2\\0\le s<a_*}}R_*^{p(r,c,s)}\\
 &\quad+\sum_{\substack{0\le r\le2,\ 0\le c\le1\\0\le s<a_*}}R_*^{H_*+p(r,c,s)}
       +R_*^{9M_*}+R_*^{27M_*}+4,\\
 M_F={}&M_F^{\rm nat}+B-1,\\
 W_0={}&R_*B^2=2^{2\cdot5^{90}+5^{45}}.
\end{align*}
The sum defining $D_C$ runs only over the nonzero $c_e$ just specified. There is no variable exponent in any final equation: these expressions define fixed integers once and for all.

\subsection{Why this is an admissible compiler slice}

We use conservative, explicit powers instead of the historical implementation's minimal radix choice. The proof of admissibility retains the actual compiler requirements. The selector radix satisfies $A_*>k_*+1$ and $A_*\equiv1\pmod5$. The mask count is
\[
 m=2\pc(\mu_*)+12a_*+2=k_*+9a_*+4,
\]
so $|\mathcal E|=m+1$ and there is exactly one dummy position. The chosen $L_*$ exceeds the highest convolution exponent plus the greatest support exponent. If $N_c=k_*+9a_*+4$, then $c_e\le14A_*^{j_*}$ and the explicit bound
\[
 G=128(m+3)N_c+16,
 \qquad b\ge100(j_*+2)+\operatorname{bitlength}(G)
\]
implies
\[
 R_*\ge \max\{4(m+3)(2\textstyle\sum_e c_e+6)+8,\ 2\mu_*+4,\ 16\}.
\]
This is the strengthened carry margin for the modified dummy bit, not merely the old Boolean-digit margin. Both $b$ and $L_*$ are powers of five.

The actual support gives $D_C\equiv4$, $D_R\equiv1\pmod5$, so $2D_C-D_R\equiv2\pmod5$. Therefore the optional high correcting term of the original compiler is unnecessary here. The masks satisfy
\[
0<M_C,M_F^{\rm nat}<B-1,\quad M_C\equiv2\pmod4,
\quad M_F^{\rm nat}\equiv4\pmod8,
\quad\pc(M_C)+\pc(M_F^{\rm nat})=d.
\]
The separate coefficient note proves these statements and the remaining clean-band and anchor requirements. The general soundness/completeness proof is the native helical compiler of raw76 and modified75, with these admissible parameters~\cite{Raw76,Raw75,HalfCompiler}. Local overlap, carry bounds, synchronization, and the positive Pell converse are all part of that dependency; arbitrary numbers with only the displayed mask congruences would not suffice.

\begin{remark}[Coefficient size is a real cost]
$B$ alone has $5^{90}+1$ binary digits. The construction is explicit as a finite mathematical definition and as a terminating algorithm, not as a practical expansion of its coefficients. The archive contains the exact definitions, geometry, machine table, and circuits. It does not contain expanded binary or decimal strings for these coefficients, and no claim of minimal coefficient height is made.
\end{remark}

\section{The seven-equation system}\label{sec:system}

The sixteen unknowns are
\[
 J,F,Z,\alpha,t,w,s,\eta,\zeta,\gamma,\tau,h,i,T,y,f\in\Pos.
\]
Define the following polynomial abbreviations:
\begin{align}
 q&=(B-1)J+1,& C&=Z+W_0,\label{eq:qC}\\
 X&=wq,&Y&=sq^3,&E&=XY,\nonumber\\
 k&=\eta+\zeta,&c&=kY+\eta,&a&=E+Y,\nonumber\\
 H&=4a+3,&\Delta&=a^2+H=(a+1)(a+3),\nonumber\\
 D_m&=X+ac+\gamma H,&L&=E(kY),\nonumber\\
 R&=\bigl(q(q-F)-Z\bigr)(q^2-1)+(M_C+qM_F)J,\label{eq:R}\\
 V&=c(Tf-1)-Rf^2,\nonumber\\
 \Lambda&=q-F-2Z-W_0-\alpha-2d.\label{eq:Lambda}
\end{align}
The fixed system $\mathcal S_{16}$ is
\begin{align}
 \Lambda&=0,\label{sys:bound}\\
 \tau^2-L(L+k)&=1,\label{sys:first}\\
 D_m^2-\Delta c^2&=1,\label{sys:main}\\
 f^2-\Delta (ic^2)^2&=1,\label{sys:strong}\\
 \Delta(f^2-1)(V^2-y^2)+y^2&=1,\label{sys:aux}\\
 k-hE-R&=1,\label{sys:index}\\
 (K_0+w)C+q-F-t(q-1)&=1.\label{sys:transport}
\end{align}
Every abbreviation is a polynomial in the sixteen displayed unknowns and fixed integers. Equation~\eqref{sys:strong} makes $\Delta(f^2-1)=(\Delta ic^2)^2$; this replacement substantially lowers the auxiliary equation's degree. Its validity is used on the strong equation, not claimed as a free polynomial identity away from that equation.

The first norm is written in a factored form:
\[
 L(L+k)=(E^2+X)(kY)^2=XY^2(XY^2+1)k^2.
\]
The identity uses the definitions $E=XY$ and $L=E(kY)$ and holds on every assignment. It saves one multiplication.

For a single equation with a straightforward proof, sum the squares of the seven residuals in~\eqref{sys:bound}--\eqref{sys:transport}. Over the integers, a sum of squares vanishes exactly when every residual vanishes. The supplied circuit for this polynomial has $86$ operations and degree $68$ in the same sixteen unknowns. Strict positivity remains a condition on these unknowns until Section~\ref{sec:domains} converts domains.

\section{Why the fixed input bridge can be deleted}\label{sec:deletion}

The repository's normalized85/84 universal construction has eighteen positive witnesses and a variable ordinary input $x$. It builds an input marker $W$ using an additional Pell norm and the three witnesses $\delta,\rho,\sigma$. Its main parameter uses $\gamma=\rho+\sigma$. For the fixed input $x=1$, the decoded marker is the known constant $W_0=2^{2d+b}$. We replace its computed word by $C=Z+W_0$, delete the three witnesses and the input norm, and supply a single independent positive $\gamma$. This gives sixteen positive witnesses.

Merely deleting an equation would be insufficient. We prove an explicit inverse that restores all omitted positive coordinates on every reduced solution. The proof applies also to the signed auxiliary-quotient domain needed later.

\subsection{The retained kernel before input decoding}

For this subsection put $A=a+2$, so $\Delta=A^2-1$ and $H=4A-5$. Define the six normalized factors
\begin{align*}
 N_1&=\tau^2-L(L+k),& N_2&=D_m^2-\Delta c^2,\\
 N_s&=f^2-\Delta i^2c^4,& Q&=(\Delta ic^2)^2,\\
 N_a&=Q(V^2-y^2)+y^2,\\
 N_k&=k-hE-R,&N_t&=(K_0+w)C+q-F-t(q-1),
\end{align*}
and $P_6=N_1N_2N_aN_kN_tN_s$. System $\mathcal S_{16}$ implies $\Lambda=0$, $P_6=1$, because its auxiliary coefficient equals $Q$ on $N_s=1$.

\begin{lemma}[Retained kernel; input witnesses not required]\label{lem:kernel}
Use the valid fixed compiler coefficients. Suppose $\Lambda=0$ and $P_6=1$, with the sixteen coordinates positive except that $T$ may be any integer, while $f>0$. Then all six factors equal $1$, $R$ is odd, and
\[
 R>3q+1,\qquad R+2<q^4\le XY<a,\qquad
 c=\psi_A(R),\quad D_m=\chi_A(R),\quad X=2^R<a.
\]
Here
\[
 \chi_A(n)+\psi_A(n)\sqrt{A^2-1}
       =(A+\sqrt{A^2-1})^n.
\]
In particular, $q$ is a power of two. No input Pell norm, split of $\gamma$, typed computation history, or positivity of $T$ is assumed.
\end{lemma}

\begin{proof}
We spell out the extraction from the existing signed-kernel proofs, because their full-parent endpoint must not be used before its missing input data are restored.

The loader equality gives $C=Z+W_0>0$, $F+Z<q$, and $W_0<q$. The native mask bounds then give
\[
 (2q-1)(q^2-1)<R<q^4-q^3,
\]
which implies the displayed pretyping inequalities. Every factor is an integer unit. The main and strong norms cannot be $-1$ modulo $4$ because $\Delta$ is $0$ or $3$ modulo $4$. The auxiliary norm is congruent to a square modulo $4$ because $Q$ is a square. The negative first norm is excluded by the standard descent for $U(U+1)$, with $U=XY^2>1$: the norm-preserving transformation
\[
 k'=(2U+1)k-2\tau,\qquad
 \tau'=(2U+1)\tau-2U(U+1)k
\]
would give $0<k'<k$ from a solution of norm $-1$, and $|\tau'|$ gives a smaller positive solution. Thus the four norm factors equal $1$, and the index and transport factors have a common sign $\epsilon\in\{\pm1\}$.

The positive first and main norms classify their pairs as
\[
 \tau=\chi_{2XY^2+1}(n),\quad k=2\psi_{2XY^2+1}(n),
 \qquad D_m=\chi_A(p),\quad c=\psi_A(p).
\]
The index congruence and ratio $kY<c<k(Y+1)$ give $p>n\ge(R-1)/2\ge24$. Strong divisibility gives $f=\chi_A(m)$, $\psi_A(m)=ic^2$, and $pc\mid m$, hence $m>2p$. These are exactly the retained rank hypotheses of the signed absorption proof~\cite{SignedAbsorption}.

The auxiliary congruences $V\equiv-c\pmod f$ and $V\equiv-R\pmod c$ hold for signed $T$. Classification of the absolute auxiliary root, followed by the strict Pell step-down at $f=\chi_A(m)$, gives $R\equiv\pm p\pmod c$. The bounds $0<R,p<c/2$ force $p=R$. This step uses $|V|$, not an unjustified assumption $V>0$.

The first-index congruence now gives $2n=R+\epsilon$. If $\epsilon=-1$, Pell duplication and $2A^2-1>2XY^2+1$ imply
\[
 c=\psi_A(2n+1)>A k>k(Y+1),
\]
contradicting the ratio interval. Thus $\epsilon=1$ and $R$ is odd.

Finally the elementary Pell ratio estimates from the signed power-recovery proof~\cite{SignedSoundness} imply, for $r=(R-1)/2$,
\[
 Y>\frac{X^r}{3},\qquad a>\frac{X^{r+1}}3>2^R.
\]
For $e_n=\chi_A(n)-a\psi_A(n)$ the recurrence gives $e_n\equiv2^n\pmod H$. The main-root definition says $e_R=X+\gamma H$. Both $X$ and $2^R$ lie strictly between $0$ and $a<H$, hence $X=2^R$. Since $q\mid X$, $q$ is dyadic.

The rank proof just used is Section~2 of signed absorption; the power proof is Section~2 of signed soundness. Their only uses of the old $\rho+\sigma$ expression are positivity of the main root and its defining projection. Both hold with our positive $\gamma$. The deleted input norm contributed only a factor already equal to $1$ to the preceding unit argument. Thus neither omitted input data nor a full-parent zero has been assumed.
\end{proof}

\subsection{A constructive, positive inverse}

Let $u=2d+b$, an odd fixed integer at least $3$. From $2^u=W_0<q\le X=2^R$ we get $u<R$. Define
\begin{align}
 \kappa&=\psi_A(u),&\mu&=\chi_A(u),\nonumber\\
 \delta&=\frac{\kappa-u}{\Delta},&
 \rho&=\frac{\mu-a\kappa-W_0}{H},&
 \sigma&=\gamma-\rho.\label{eq:restore}
\end{align}
These quotients occur in the inverse proof, not in the polynomial circuit.

For odd $u$, the Pell recurrence modulo $\Delta$ gives $\psi_A(u)\equiv u\pmod\Delta$. Since $u\ge3$, $\psi_A(u)>u$, so $\delta$ is a positive integer. For $e_n=\chi_A(n)-a\psi_A(n)$,
\[
 e_0=1,\quad e_1=2,\quad e_{n+1}=2Ae_n-e_{n-1}.
\]
As $H=4A-5$, the sequence $2^n$ satisfies this recurrence modulo $H$. Thus $\rho$ is integral. Strict growth gives $e_{n+1}>(2A-1)e_n>2e_n$ for $n\ge1$, so $e_u>2^u$ and $\rho>0$.

The remaining positivity is not automatic and is the essential check. Since $R>u\ge1$,
\[
 e_R-e_u\ge e_{u+1}-e_u>(2A-2)e_u\ge4A-4>A>X=2^R.
\]
The main-root equation therefore implies
\[
 H\sigma=e_R-e_u-X+2^u>0.
\]
So $\sigma$ is also a positive integer. The identities
\[
 \kappa=u+\delta\Delta,\qquad
 \mu=W_0+a\kappa+\rho H,\qquad
 \mu^2-\Delta\kappa^2=1,\qquad \gamma=\rho+\sigma
\]
restore exactly the omitted input cone of the normalized85/84 source at $x=1$. If desired, the earlier positive gap is $\psi_A(R)-\psi_A(u)>0$.

This yields a full parent zero with every supplied witness positive, except possibly the already allowed signed $T$. For positive $T$, the ordinary compiler theorem now applies. Conversely, a positive full-parent zero at $x=1$ has $W=W_0$ by its existing input theorem. Set $\gamma=\rho+\sigma$ and delete $\delta,\rho,\sigma$; all seven equations of $\mathcal S_{16}$ hold. This proves exact existence equivalence with the fixed-input parent, not only a map on canonical padding choices.

\section{Two explicit polynomials in fifteen unknowns}\label{sec:fifteen}

Delete $f$ from the sixteen-variable list. Keep all definitions not involving $f$. Define
\begin{align}
 Q&=(\Delta ic^2)^2,&\mathcal D&=1+\Delta i^2c^4,\nonumber\\
 v&=cT,&z&=c+R\mathcal D,\nonumber\\
 \mathcal A&=Q(\mathcal D v^2+z^2-y^2)+y^2-1,\label{eq:rootdata}\\
 \mathcal B&=2Qvz,&\mathcal K&=\mathcal A^2-\mathcal D\mathcal B^2,\nonumber\\
 P_4&=N_1N_2N_kN_t.\nonumber
\end{align}
All these expressions are integer polynomials in fifteen unknowns. Our principal single equation is
\begin{equation}\boxed{\quad\Phi_{15}=P_4(1+\Lambda^2)(1+\mathcal K)-1=0.\quad}\label{eq:Phi}
\end{equation}
The circuit computes $\mathcal D\mathcal B^2$ as
\[
4(Q\mathcal D v^2)(Qz^2),
\]
so it need not separately construct $\mathcal B$. There is no uncounted division or existential root in~\eqref{eq:Phi}.

\subsection{Soundness, including the sign of the restored root}

At an integer zero of~\eqref{eq:Phi}, every displayed factor is a unit. Since $1+\Lambda^2\ge1$, it must equal $1$, giving $\Lambda=0$. Also $\mathcal B$ is even, so
\[
 1+\mathcal K\equiv1+\mathcal A^2\pmod4
\]
is $1$ or $2$ modulo $4$ and cannot be $-1$. Hence $\mathcal K=0$ and $P_4=1$.

The loader bounds imply $R>0$; therefore $z,v,Q,\mathcal B>0$. In fact nonvanishing also follows before those bounds from $\mathcal D>c>0$, since $c+R\mathcal D$ cannot be zero for integral $R$. Put
\[
 f=\mathcal A/\mathcal B\in\QQ.
\]
Then $f^2=\mathcal D\in\Pos$. A rational number with integral square is integral: in a reduced fraction $p/q$, the divisibility $q^2\mid p^2$ forces $q=1$. Thus $f$ is a nonzero integer. On $f^2=\mathcal D$ we have $V=vf-z$, and direct expansion gives
\[
 QV^2-(Q-1)y^2=\mathcal A+1-\mathcal Bf=1,
 \qquad f^2-\Delta i^2c^4=1.
\]
These restore the six-factor system, provisionally with signed $f$.

To obtain positivity without circular reasoning, replace $(f,T)$ by
\[
 (f_+,T_+)=(|f|,\operatorname{sign}(f)T).
\]
Every retained expression is unchanged, because the only occurrences are $f^2$ and $Tf$. Lemma~\ref{lem:kernel} applies with $f_+>0$ and signed $T_+$. The explicit inverse~\eqref{eq:restore} then supplies positive $\delta,\rho,\sigma$. We have now obtained a \emph{full} parent84 zero with only its auxiliary quotient allowed signed.

The established signed-quotient theorem~\cite{SignedSoundness} applies to precisely that domain and proves $T_+>0$. Its proof uses odd-index Pell ratios, the exact population bound, the actual compiler masks to recover $R\equiv3\pmod4$, and both unsquared auxiliary congruences to fix the sign. Since the original $T>0$, this forces $f>0$. Thus every positive zero of $\Phi_{15}$ restores a positive solution of $\mathcal S_{16}$.

The order matters: applying the signed-quotient theorem before restoring its input tuple, or applying the positive compiler while $f$ is still signed, would be invalid. The retained-kernel extraction and explicit inverse supply the missing steps.

Conversely, at a positive solution of $\mathcal S_{16}$, the strong equation gives $f^2=\mathcal D$ and the auxiliary equation gives $\mathcal A=\mathcal Bf$. Thus $\mathcal K=0$, $P_4=1$, and $\Lambda=0$. Deleting $f$ gives a positive zero of $\Phi_{15}$. The restored $f$ is unique because $\mathcal B\ne0$.

\subsection{An eighty-operation alternative}

Put $P=P_4(1+\Lambda^2)$. A second fifteen-variable polynomial is
\begin{equation}\boxed{\quad\Psi_{15}=[P(\mathcal A+1)-1]^2-\mathcal D(P\mathcal B)^2=0.\quad}\label{eq:Psi}
\end{equation}
Its supplied circuit shares the two large products rather than evaluating this display literally.

Here is a direct proof of equivalence. At a zero, expansion shows $1$ is a multiple of $P$, hence $P=\pm1$. This forces $\Lambda=0$. Define
\[
 f=\frac{\mathcal A+1-P}{\mathcal B}.
\]
Equation~\eqref{eq:Psi} gives $f^2=\mathcal D$, so $f$ is a nonzero integer. Its auxiliary norm is
\[
 QV^2-(Q-1)y^2=\mathcal A+1-\mathcal Bf=P.
\]
Because $Q$ is a square, this expression is a square modulo $4$ and cannot equal $-1$. Therefore $P=1$ and $f=\mathcal A/\mathcal B$, yielding $\mathcal K=0$ and a zero of $\Phi_{15}$. The converse is immediate. This proof needs no unproved sign choice for a rational root; the same signed-root argument above finishes the positive-domain equivalence.

\section{Degrees and arithmetic operation counts}

The archive contains every arithmetic row, every free port, and every final comparison or output. The names in the mathematical display map to source registers in the supplied documentation. Each circuit is acyclic; all its emitted rows and fixed ports are live. Its computed registers are signed polynomial values, not additional positive unknowns.

Let $q_0=(B-1)J$ and $k_0=\eta+\zeta$. The main norm's nominal highest terms cancel:
\[
 (X+ac+g)^2-(a^2+H)c^2
 =X^2+2Xg+g^2+2acX+2acg-Hc^2,\quad g=\gamma H.
\]
Its exact degree is $18$, with leading form $8\gamma k_0w^2s^3q_0^{11}$. The first norm has degree $22$; the index and transport factors have degrees $7$ and $2$. Thus $P_4$ has degree $49$. Its leading form is
\[
 8\gamma h k_0^3w^5s^8q_0^{29}(wZ-tq_0),
\]
a nonzero polynomial for every admissible $B>1$.

In the reduced auxiliary system, the seven equation degrees are, in the order~\eqref{sys:bound}--\eqref{sys:transport},
\[
 1,\quad22,\quad18,\quad34,\quad28,\quad7,\quad2.
\]
The degree-$34$ strong residual has leading form
\[
 -i^2k_0^4w^2s^6q_0^{20}.
\]
Its square is the unique degree-$68$ summand in the sum-of-squares polynomial, proving exact degree $68$ uniformly.

For the fifteen-variable construction,
\[
 \deg Q=46,\quad\deg\mathcal D=34,\quad\deg z=38,
 \quad\deg\mathcal A=122,\quad\deg\mathcal B=90.
\]
The leading form of $\mathcal A$ is $Q_{\rm top}R_{\rm top}^2\mathcal D_{\rm top}^2$. Thus $\mathcal K$ has exact degree $244$: its square term has degree $244$, while $\mathcal D\mathcal B^2$ has degree $214$. Since $\Lambda$ has the nonzero degree-one part $q_0-F-2Z-\alpha$,
\[
 \deg\Phi_{15}=49+2+244=295.
\]
For $\Psi_{15}$, $\deg P=51$. The first squared term has degree $2(51+122)=346$ and its subtracted term has degree $34+2(51+90)=316$, giving exact degree $346$. These are uniform degree arguments, not conclusions from random evaluation or zero-locus simplification.

The principal ledgers are
\begin{center}
\begin{tabular}{lrrr}
\toprule
Circuit & Multiplications & Additions/subtractions & Total\\
\midrule
Seven-equation reduced system &35&31&66\\
Its complete sum of squares &42&44&86\\
Fifteen-variable local norm guard &46&36&82\\
Fifteen-variable outer norm &46&34&80\\
\bottomrule
\end{tabular}
\end{center}
The system's $66$ operations are shared across its seven comparisons. Its sum-of-squares finalizer adds seven residual subtractions, seven squares, and six additions. Other sixteen-variable product tradeoffs, with their exact ledgers and degrees, are included in the source catalogue; they use the same positive zero set and no extra theorem about K215.

\section{Nonnegative and unrestricted integer solutions}\label{sec:domains}

\subsection{Nonnegative integers with the same number of unknowns}

For any primary polynomial $P(z_1,\ldots,z_n)$, define
\[
 P_{\NN}(x_1,\ldots,x_n)=P(x_1+1,\ldots,x_n+1).
\]
Then $P_{\NN}=0$ has a solution in $\NN^n$ exactly when $P=0$ has a solution in $\Pos^n$. Apply the same substitution to each equation of a system. Degrees are unchanged; the direct shifted circuit adds $n$ additions unless further simplification is performed. The displayed primary gate counts concern the positive-variable circuits.

\subsection{A single equation over all integers}

Use the integer polynomial
\[
 \theta(u,v,w)=1+u^2+v^2+w(w+1).
\]
It is positive at every integer triple and takes every positive integer value. To prove the latter claim, apply Legendre's three-square theorem to $4(z-1)+1$: this integer is neither divisible by four nor congruent to seven modulo eight. Exactly one of its three square summands is odd and the other two are even. Writing them as $(2w+1)^2,(2u)^2,(2v)^2$ gives $z=\theta(u,v,w)$. This standard reduction is also used in modern universal-pair formalizations~\cite{BayerDavid2025}.

If $(z_1,\ldots,z_{15})$ lists the fifteen positive unknowns, define the completely unrestricted polynomial
\begin{equation}
 \widehat\Phi_{45}(u_1,v_1,w_1,\ldots,u_{15},v_{15},w_{15})
 =\Phi_{15}\bigl(\theta(u_1,v_1,w_1),\ldots,\theta(u_{15},v_{15},w_{15})\bigr).
\end{equation}
It has integer coefficients, exact degree $590$, and a zero in $\ZZ^{45}$ exactly when $\Phi_{15}$ has a positive zero. Hence
\[
 \mathrm{ZFC}\vdash\Con(\mathrm{ZFC})\ \longleftrightarrow\
 \forall\mathbf u\in\ZZ^{45}\;\widehat\Phi_{45}(\mathbf u)\ne0.
\]
The exact degree doubles: the substituted top form remains nonzero, as can be seen by setting every $v_j,w_j$ to zero and retaining the independent $u_j^2$.

The archive also emits the same substitution for the sixteen-variable sum of squares, giving one equation in $48$ unrestricted integer unknowns of exact degree $136$. A direct evaluation of $\theta$ costs three multiplications and four additions. Therefore the complete $45$-variable circuit costs $187=91\mathrm M+96\mathrm A$ operations, and the complete $48$-variable circuit costs $198=90\mathrm M+108\mathrm A$. Their emitted JSON circuits include every substitution operation. Lagrange's familiar four-square conversion would instead require $60$ or $64$ unrestricted unknowns.

Allowing the primary variables to range over arbitrary signed integers without a conversion is not valid. In particular, scaled forms can acquire spurious zeros where $\Delta=0$. The domain is part of the theorem, not a numerical implementation detail.

\section{Why the equivalence is provable in ZFC}

Let $\operatorname{Ref}_{\mathrm{ZFC}}$ denote the existence of a finite derivation of falsity in the fixed ordinary ZFC calculus, and define $\Con(\mathrm{ZFC})=\neg\operatorname{Ref}_{\mathrm{ZFC}}$. The proof establishes the following chain by explicit finite transformations and arithmetic lemmas:
\[
\begin{aligned}
 \Sol(\Phi_{15})
 &\longleftrightarrow\Sol(\mathcal S_{16})\\
 &\longleftrightarrow\text{the normalized machine accepts four unary ones}\\
 &\longleftrightarrow\operatorname{Halt}(K215,\mathrm{blank})\\
 &\longleftrightarrow\operatorname{Ref}_{\mathrm{ZFC}}.
\end{aligned}
\]
The same holds for $\Psi_{15}$ and the stated domain conversions. Negating both sides gives Theorem~\ref{thm:main}.

Each new argument is ordinary mathematics about integers, finite strings, finite tables, finite computations, recurrence sequences, divisibility, and polynomial identities. ZFC proves the existence and elementary properties of these objects and carries out the indicated inductions. The inherited compiler proofs and K215 calibration are invoked as their ordinary mathematical theorems. We do not infer an internal ZFC theorem merely from the fact that Lean or Rocq accepted an external statement, and we do not claim that ZFC interprets every universe of either assistant. The accepted candidate's mathematical calibration already states the ZFC-provable endpoint; the external kernel checks provide separate evidence for its formalized counterpart.

In particular, no consistency assumption is used to prove the equivalence. The equivalence does not prove its own no-solution side. If ZFC is consistent and has the usual effective arithmetic presentation, a ZFC proof of that no-solution assertion would yield a ZFC proof of $\Con(\mathrm{ZFC})$, contrary to the second incompleteness theorem. Consistency alone also should not be described as ruling out every ZFC proof of the opposite $\Sigma_1$ assertion; stronger soundness assumptions are needed for that separate claim.

\section{Reproducible evidence and remaining formalization work}\label{sec:evidence}

The archive supplies the paper source and PDF, the literal dense and native K215 tables and state map, the normalized finite machine, the exact finite coefficient specification, seven positive-domain circuit variants and two unrestricted-integer lifts, standalone circuit builders, independent expression-transcription checks, and a top-level verification command. Its provenance manifest pins the inherited bytes and new outputs. New source files follow the repository's MIT No Attribution terms; the inherited machine notices and third-party license notices are retained separately.

The checks establish the following finite and algebraic facts:
\begin{itemize}[nosep]
\item exact source hashes and all $430$ dense/native transition correspondences;
\item every folded native instruction context, including origin crossing and both tracks;
\item the finite normalized table, forced first-step convention, Start/End membership, and factored window count;
\item the exact coefficient geometry, strengthened radix bounds, mask-population formulas, and required five-adic unit;
\item literal circuit reconstruction, supplied-variable counts, topological order, liveness, and complete operation ledgers;
\item independent signed and rational evaluations of every full emitted circuit against its mathematical expression;
\item exact local polynomial identities, plus full univariate coefficient diagnostics supporting the separate uniform degree proofs.
\end{itemize}
The tests do not enumerate all accepting histories, expand the enormous fixed coefficient integers, construct a full Pell witness, or decide whether K215 halts. Finite tests supplement the general proof; they are not substituted for universal soundness or completeness. In particular, the historical H10 compiler and the signed quotient theorem remain mathematical proof dependencies, not freshly compiled Lean/Rocq theorems in this packet.

For integration, a useful next formal milestone is a compiler theorem with a fixed polynomial expression datatype and a literal statement that its positive zero set is equivalent to acceptance of the normalized finite machine. The existing K215 endpoint then gives the last logical composition immediately. A theorem with that compiler equivalence left as an explicit hypothesis would be a conditional wrapper, not a verification of the new equation. We do not count such a wrapper as completed formal work.

\section{Further research questions}

\begin{enumerate}[leftmargin=*,itemsep=4pt]
\item \textbf{Compile the numerical calculus directly.} The repository already proves a twelve-rule numerical calculus equivalent to ZF refutability, and its constructor codes use simple quadratic expressions. Bypassing the two-way TM and the very large window alphabet could reduce the coefficient description and height substantially, even if the arithmetic circuit's variable count stays unchanged.
\item \textbf{Reduce the fixed-input slack.} Our loader retains both the positive field witness $F$ and the positive upper-bound slack $\alpha$. Can one be removed while proving the pretyping range needed by the native compiler? Dropping it without a replacement range proof is unsound.
\item \textbf{Shorten the signed-root proof.} The fifteen-variable construction depends on the signed auxiliary quotient theorem. A direct positivity criterion expressed in the retained norms could simplify the proof substantially or produce a different fifteen-variable equation with lower degree.
\item \textbf{Use fixed input earlier.} We remove the input Pell bridge after the universal polynomial is built. A tableau compiler with a single fixed initial pattern might eliminate the marker mechanism itself and reduce the number of arithmetic witnesses further.
\item \textbf{Optimize several costs together.} The $(80,346,15)$ and $(82,295,15)$ points trade operations against degree. Search for improvements with an explicit cost tuple including coefficient-description length and bit height, not only gate count.
\item \textbf{Formalize the bridge in small modules.} Separate finite machine normalization, admissible window compilation, Pell/population recovery, fixed-input deletion, and root elimination. The local algebra and positive-inverse lemmas should be substantially smaller than the inherited full-machine endpoint.
\end{enumerate}

\appendix
\section{Complete circuit catalogue and source names}

\begin{center}\small
\begin{tabular}{lrrrr}
\toprule
JSON stem & Variables & Equations & Degree & Gates\\
\midrule
\texttt{fixed16\_system66\_degree34} &16&7&34&66\\
\texttt{fixed16\_sos86\_degree68} &16&1&68&86\\
\texttt{fixed16\_product74\_degree157} &16&1&157&74\\
\texttt{fixed16\_product75\_degree125} &16&1&125&75\\
\texttt{fixed16\_product76\_degree113} &16&1&113&76\\
\texttt{fixed15\_product82\_degree295} &15&1&295&82\\
\texttt{fixed15\_outer80\_degree346} &15&1&346&80\\
\midrule
\texttt{integer45\_product187\_degree590} &45&1&590&187\\
\texttt{integer48\_sos198\_degree136} &48&1&136&198\\
\bottomrule
\end{tabular}\end{center}

The first seven rows use strictly positive unknowns. The final two rows allow all integer values and include every operation of the three-variable substitution in Section~\ref{sec:domains}.

The sixteen-variable product forms use the six normalized factors, or their scaled strong version, multiplied by $1+\Lambda^2$. The $74$-gate version uses the shared square coefficient $Q$ and the scaled strong factor $\Delta N_s$, subtracting $\Delta$ at the output. The $75$-gate version replaces its auxiliary coefficient by $\Delta(f^2-1)$. The $76$-gate version uses the latter coefficient and the unscaled strong factor, subtracting $1$ at the output. At any zero, the strong factor is a unit and cannot be $-1$ modulo $4$, so it is $1$ and both auxiliary coefficients agree. The remaining factor recovery is Lemma~\ref{lem:kernel}. Conversely each system solution clearly gives each product zero. This proves their common positive solution condition without treating equality on the strong equation as an unconditional polynomial identity.

\begin{center}\small
\begin{tabular}{ll@{\qquad}ll}
\toprule
Mathematics & Source name & Mathematics & Source name\\
\midrule
$J$&\texttt{Jrep}&$t$&\texttt{transport\_quotient}\\
$T$&\texttt{auxiliary\_quotient}&$\tau$&\texttt{tau\_root}\\
$y$&\texttt{y\_aux}&$\Delta$&\texttt{A}\\
$c$&\texttt{R10a}&$a$&\texttt{R12}\\
$X$&\texttt{wn2}&$Y$&\texttt{sn2}\\
$E$&\texttt{UM}&$R$&\texttt{r\_lhs}\\
$C$&\texttt{marked\_rhs}&$K_0$&\texttt{Kconstant}\\
\bottomrule
\end{tabular}\end{center}
The source name \texttt{A} means $\Delta$, not the Pell parameter $a+2$. The fixed source name \texttt{MF} means $M_F^{\rm nat}+B-1$.

\section{Inherited mathematical interfaces}

The retained-kernel proof uses the signed absorption note, Section~2, and the signed soundness note, Section~2, before any full-parent theorem. Positive input restoration then permits invocation of the full normalized85/84 theorem. Fifteen-variable sign restoration additionally uses the full signed soundness theorem, after the omitted input tuple has been restored. The machine/compiler path uses raw81's stationary first-step convention, raw76's Start/End ordering and asymmetric temporal scale, and modified75's masks and half-binomial population rule. The archive's provenance file records exact paths and hashes rather than mutable branch links.

\begin{thebibliography}{99}\small
\bibitem{K215Verification} ProveIt contributors, \emph{Formal correctness of the exact 215-state machine}, pinned repository verification page.
\href{https://github.com/VladimirReshetnikov/ProveIt/blob/cdd582bec13f908df3711fa9e05d1ddac368a5db/Computability/BusyBeaver/Research/ZFC/k225-formal/verification/k215/README.md}{Immutable source.}
\bibitem{K215Candidate} ProveIt contributors, \emph{K215: 215 working states for ZFC consistency}, accepted candidate, proof exposition, and notices.
\href{https://github.com/VladimirReshetnikov/ProveIt/tree/cdd582bec13f908df3711fa9e05d1ddac368a5db/Computability/BusyBeaver/Research/ZFC/candidate-k215}{Immutable source directory.}
\bibitem{Raw76} ProveIt contributors, \emph{A 76-operation fixed-index universal certificate}, with the preceding raw81 and helical unary-tableau constructions. \href{https://github.com/VladimirReshetnikov/ProveIt/blob/cdd582bec13f908df3711fa9e05d1ddac368a5db/Computability/HilbertTenthProblem/Papers/1980/FIXED_RAW_UNIVERSAL_76_PROOF.md}{Immutable proof.}
\bibitem{Raw75} ProveIt contributors, \emph{A 75-operation fixed-index universal certificate}. \href{https://github.com/VladimirReshetnikov/ProveIt/blob/cdd582bec13f908df3711fa9e05d1ddac368a5db/Computability/HilbertTenthProblem/Papers/1980/FIXED_RAW_UNIVERSAL_75_PROOF.md}{Immutable proof.}
\bibitem{HalfCompiler} ProveIt contributors, \emph{Modified native masks for the half-binomial complete75 compiler}. \href{https://github.com/VladimirReshetnikov/ProveIt/blob/cdd582bec13f908df3711fa9e05d1ddac368a5db/Computability/HilbertTenthProblem/Papers/research-wip/native-stream-queue/complete75_half_binomial_compiler.md}{Immutable proof and source.}
\bibitem{Normalized85} ProveIt contributors, \emph{An auxiliary quotient gives a complete universal polynomial in 85 operations}. \href{https://github.com/VladimirReshetnikov/ProveIt/blob/cdd582bec13f908df3711fa9e05d1ddac368a5db/Computability/HilbertTenthProblem/Papers/research-wip/native-stream-queue/complete85_auxiliary_bezout_projection.md}{Immutable proof.}
\bibitem{Scaled84} ProveIt contributors, \emph{A complete universal polynomial in 84 operations}. \href{https://github.com/VladimirReshetnikov/ProveIt/blob/cdd582bec13f908df3711fa9e05d1ddac368a5db/Computability/HilbertTenthProblem/Papers/research-wip/native-stream-queue/complete84_scaled_strong_output.md}{Immutable proof and circuit.}
\bibitem{SignedAbsorption} ProveIt contributors, \emph{Signed auxiliary quotient absorption}, especially the signed pre-input rank argument in Section~2. \href{https://github.com/VladimirReshetnikov/ProveIt/blob/cdd582bec13f908df3711fa9e05d1ddac368a5db/Computability/HilbertTenthProblem/Papers/research-wip/native-stream-queue/complete84_signed_quotient_absorption.md}{Immutable proof.}
\bibitem{SignedSoundness} ProveIt contributors, \emph{Signed auxiliary quotient: source-specific recovery of the positive domain}, and its independent mathematical review. \href{https://github.com/VladimirReshetnikov/ProveIt/blob/cdd582bec13f908df3711fa9e05d1ddac368a5db/Computability/HilbertTenthProblem/Papers/research-wip/native-stream-queue/complete84_signed_quotient_soundness_scout.md}{Immutable proof.}
\bibitem{Norm92} ProveIt contributors, \emph{A 92-operation universal polynomial with 17 positive witnesses and degree 325}. \href{https://github.com/VladimirReshetnikov/ProveIt/blob/cdd582bec13f908df3711fa9e05d1ddac368a5db/Computability/HilbertTenthProblem/Papers/research-wip/native-stream-queue/complete92_norm_unit_guard.md}{Immutable norm-elimination proof.}
\bibitem{Jones1982} J.~P. Jones, \emph{Universal Diophantine equation}, Journal of Symbolic Logic \textbf{47} (1982), 549--571. \href{https://doi.org/10.2307/2273588}{doi:10.2307/2273588}.
\bibitem{BayerDavid2025} J.~Bayer and M.~David, \emph{A Formal Proof of Complexity Bounds on Diophantine Equations}, ITP 2025, LIPIcs \textbf{352}, article 3, especially Sections 1.2 and 3.2. \href{https://doi.org/10.4230/LIPIcs.ITP.2025.3}{doi:10.4230/LIPIcs.ITP.2025.3}; \href{https://arxiv.org/html/2505.16963v2}{arXiv:2505.16963v2}.
\end{thebibliography}
\end{document}
